\documentclass[11pt]{article}

\usepackage{amsmath,amssymb,amsthm,mathtools}
\usepackage[margin=1in]{geometry}
\usepackage{microtype}
\usepackage{enumitem}

\newtheorem{theorem}{Theorem}[section]
\newtheorem{lemma}[theorem]{Lemma}
\newtheorem{proposition}[theorem]{Proposition}
\newtheorem{corollary}[theorem]{Corollary}
\newtheorem{remark}[theorem]{Remark}
\newtheorem{definition}[theorem]{Definition}

\newcommand{\R}{\mathbb{R}}
\newcommand{\E}{\mathbb{E}}
\newcommand{\1}{\mathbf{1}}
\newcommand{\PhiH}{\Phi}
\newcommand{\Hom}{\operatorname{Hom}}
\newcommand{\tr}{\operatorname{Tr}}
\newcommand{\ip}[2]{\left\langle #1,#2\right\rangle}
\newcommand{\dd}{\,\mathrm{d}}
\newcommand{\join}{\vee}

\title{A Universal-Vertex Lifting Theorem for Goodman-Type Graphon Bounds\\
and an Application to Clique Joins of Odd Cycles}
\author{}
\date{}

\begin{document}

\maketitle

\begin{abstract}
For a finite graph $H$, define
\[
\Phi_H(p)
:=
(1-p)^{v(H)}
\chi_H\!\left(\frac{1}{1-p}\right),
\]
with the value at $p=1$ defined by continuous extension.
We prove a general lifting theorem showing that, under a simple convexity
hypothesis on $\Phi_F$, a Goodman-type lower bound for a graph $F$
automatically yields the corresponding lower bound for $K_1\join F$.
The crucial input is a weighted local Goodman inequality controlling
the average edge density in the graphon links.

Starting from the previously established odd-cycle inequality
\[
t(C_m,W)\ge p^m-p(1-p)^{m-1},
\qquad m\ge 3\text{ odd},
\]
the lifting theorem can be iterated.  As a consequence, for every
$s\ge 0$ and every odd $m\ge 3$,
\[
H_{s,m}:=K_s\join C_m
\]
satisfies
\[
t(H_{s,m},W)\ge \Phi_{H_{s,m}}(p)
\]
whenever
\[
p\ge 1-\frac1{s+2}.
\]
Since $\chi(H_{s,m})=s+3$, this is stronger than the relaxed range
$p\ge 1-1/\chi(H_{s,m})$.
\end{abstract}


%%%%%%%%%%%%%%%%%%%%%%%%%%%%%%%%%%%%%%%%%%%%%%%%%%%%%%%%%%%%%%%%%%%%%%%%%%%%%%%
\section{Setup}
%%%%%%%%%%%%%%%%%%%%%%%%%%%%%%%%%%%%%%%%%%%%%%%%%%%%%%%%%%%%%%%%%%%%%%%%%%%%%%%

We work on an atomless standard probability space
$(\Omega,\mu)$.
A graphon is a symmetric measurable function
\[
W:\Omega^2\to[0,1].
\]
For a finite graph $H$,
\[
t(H,W)
=
\int_{\Omega^{V(H)}}
\prod_{uv\in E(H)}
W(x_u,x_v)
\prod_{u\in V(H)}\dd\mu(x_u).
\]
The edge density of $W$ is
\[
p=t(K_2,W)
=
\int_{\Omega^2}W(x,y)\,\dd\mu(x)\dd\mu(y).
\]

For a finite graph $H$, let
\[
\Phi_H(p)
=
(1-p)^{v(H)}
\chi_H\!\left(\frac1{1-p}\right),
\qquad 0\le p<1.
\]
Although this expression appears singular at $p=1$, it is in fact a
polynomial in $1-p$.  Indeed, if
\[
\chi_H(x)=\sum_{j=0}^{v(H)}c_jx^j,
\]
then
\[
(1-p)^{v(H)}
\chi_H\!\left(\frac1{1-p}\right)
=
\sum_{j=0}^{v(H)}
c_j(1-p)^{v(H)-j}.
\]
We therefore define $\Phi_H(1)$ by continuous extension.
Since $\chi_H$ is monic, $\Phi_H(1)=1$.

The graph join $G\join H$ is obtained from the disjoint union
of $G$ and $H$ by adding every edge between $V(G)$ and $V(H)$.

We shall use the following result as an established input.

\begin{theorem}[Odd-cycle Goodman bound: standing input]
\label{thm:odd-cycle-input}
Let $m\ge 3$ be odd.  For every graphon $W$ with edge density
$p=t(K_2,W)$,
\[
t(C_m,W)
\ge
g_m(p),
\]
where
\[
g_m(p)
:=
p^m-p(1-p)^{m-1}.
\]
Equivalently,
\[
t(C_m,W)\ge \Phi_{C_m}(p).
\]
\end{theorem}

The purpose of this note is not to reprove
Theorem~\ref{thm:odd-cycle-input}, but to derive a substantially larger
family of graphs from it.


%%%%%%%%%%%%%%%%%%%%%%%%%%%%%%%%%%%%%%%%%%%%%%%%%%%%%%%%%%%%%%%%%%%%%%%%%%%%%%%
\section{Links of a vertex in a graphon}
%%%%%%%%%%%%%%%%%%%%%%%%%%%%%%%%%%%%%%%%%%%%%%%%%%%%%%%%%%%%%%%%%%%%%%%%%%%%%%%

For a graphon $W$, define its degree function
\[
d(x):=\int_\Omega W(x,y)\,\dd\mu(y).
\]
Then
\[
\int_\Omega d(x)\,\dd\mu(x)=p.
\]

For each $x\in\Omega$ with $d(x)>0$, define a probability measure
$\mu_x$ on $\Omega$ by
\[
\dd\mu_x(y)
=
\frac{W(x,y)}{d(x)}\,\dd\mu(y).
\]
Because $\mu_x$ is absolutely continuous with respect to the atomless
measure $\mu$, it is itself atomless.  Thus $W$, viewed as a kernel on
the probability space $(\Omega,\mu_x)$, is again a graphon.
We denote this graphon by $W_x$.

Define also the rooted triangle density
\[
\tau(x)
:=
\int_{\Omega^2}
W(x,y)W(x,z)W(y,z)
\,\dd\mu(y)\dd\mu(z).
\]
The edge density of $W_x$ is
\[
z_x
:=
t(K_2,W_x)
=
\frac{\tau(x)}{d(x)^2}
\qquad
\text{whenever }d(x)>0.
\]
When $d(x)=0$, we define $z_x:=0$; this value is immaterial in all
subsequent expressions because it is multiplied by a positive power
of $d(x)$.

The following identity is the basic reason that universal vertices
are naturally related to graphon links.

\begin{lemma}[Universal-vertex link identity]
\label{lem:link-identity}
Let $F$ be a graph on $n$ vertices and let
\[
F^+:=K_1\join F.
\]
Then
\[
t(F^+,W)
=
\int_\Omega
d(x)^n\,t(F,W_x)\,\dd\mu(x).
\]
\end{lemma}

\begin{proof}
Let the universal vertex of $F^+$ be denoted by $0$, and write
$V(F)=\{1,\dots,n\}$.  By definition,
\begin{align*}
t(F^+,W)
&=
\int_{\Omega^{n+1}}
\left(
\prod_{i=1}^n W(x_0,x_i)
\right)
\left(
\prod_{ij\in E(F)}W(x_i,x_j)
\right)
\prod_{i=0}^n\dd\mu(x_i).
\end{align*}
Fix $x_0=x$.  If $d(x)>0$, then
\[
W(x,x_i)\,\dd\mu(x_i)
=
d(x)\,\dd\mu_x(x_i).
\]
Therefore the inner integral over $x_1,\dots,x_n$ is
\[
d(x)^n
\int_{\Omega^n}
\prod_{ij\in E(F)}W(x_i,x_j)
\prod_{i=1}^n\dd\mu_x(x_i)
=
d(x)^n t(F,W_x).
\]
If $d(x)=0$, both sides of this identity vanish.  Integrating in $x$
proves the claim.
\end{proof}


%%%%%%%%%%%%%%%%%%%%%%%%%%%%%%%%%%%%%%%%%%%%%%%%%%%%%%%%%%%%%%%%%%%%%%%%%%%%%%%
\section{A weighted local Goodman inequality}
%%%%%%%%%%%%%%%%%%%%%%%%%%%%%%%%%%%%%%%%%%%%%%%%%%%%%%%%%%%%%%%%%%%%%%%%%%%%%%%

The next estimate is the key analytic input.

\begin{lemma}[Weighted rooted-triangle inequality]
\label{lem:weighted-local-goodman}
Let $W$ be a graphon with edge density $p>0$ and degree function $d$.
For every integer $k\ge 0$,
\[
\boxed{
\int_\Omega d(x)^k\tau(x)\,\dd\mu(x)
\ge
\frac{2p-1}{p}
\int_\Omega d(x)^{k+2}\,\dd\mu(x).
}
\]
\end{lemma}

\begin{proof}
For $j\ge0$, write
\[
M_j:=\int_\Omega d(x)^j\,\dd\mu(x).
\]
In particular,
\[
M_1=p.
\]

Let $T_W$ denote the graphon operator
\[
(T_Wf)(x)
=
\int_\Omega W(x,y)f(y)\,\dd\mu(y).
\]
Set
\[
I_k
:=
\int_\Omega d(x)^k\tau(x)\,\dd\mu(x),
\]
and
\[
J_k
:=
\int_{\Omega^2}
d(x)^kW(x,y)d(y)
\,\dd\mu(x)\dd\mu(y)
=
\ip{d^k}{T_Wd}.
\]

For $a,b\in[0,1]$,
\[
ab\ge a+b-1.
\]
Applying this with
\[
a=W(x,y),\qquad b=W(x,z),
\]
we obtain
\[
W(x,y)W(x,z)
\ge
W(x,y)+W(x,z)-1.
\]
Multiplying by
\[
d(x)^kW(y,z)\ge0
\]
and integrating over $x,y,z$ gives
\[
I_k\ge 2J_k-pM_k.
\tag{3.1}
\label{eq:Ik-Jk}
\]

We next estimate $J_k$.

Put
\[
U:=1-W
\]
and let $T_U$ be its graphon operator.  Define
\[
u:=T_U\1.
\]
Since
\[
d=T_W\1
\]
and $W+U=1$, we have
\[
u=1-d.
\]
Also
\[
\int_\Omega u\,\dd\mu
=
1-p.
\]
Write
\[
q:=1-p.
\]

Since
\[
T_Ud=T_U(\1-u)=u-T_Uu,
\]
we have
\[
T_Wd
=
\left(\int_\Omega d\,\dd\mu\right)\1-T_Ud
=
p\1-u+T_Uu
=
d-q\1+T_Uu.
\]
Consequently,
\begin{equation}
J_k
=
M_{k+1}-qM_k+\ip{d^k}{T_Uu}.
\tag{3.2}
\label{eq:Jk-decomp}
\end{equation}

We claim that
\begin{equation}
\ip{d^k}{T_Uu}
\ge
\int_\Omega d(x)^ku(x)^2\,\dd\mu(x).
\tag{3.3}
\label{eq:correlation}
\end{equation}

Indeed, by symmetry of $U$,
\begin{align*}
2\ip{d^k}{T_Uu}
&=
\int_{\Omega^2}
U(x,y)
\left(
d(x)^ku(y)+d(y)^ku(x)
\right)
\,\dd\mu(x)\dd\mu(y),
\end{align*}
whereas, since
\[
u(x)=\int_\Omega U(x,y)\,\dd\mu(y),
\]
we have
\begin{align*}
2\int_\Omega d(x)^ku(x)^2\,\dd\mu(x)
&=
\int_{\Omega^2}
U(x,y)
\left(
d(x)^ku(x)+d(y)^ku(y)
\right)
\,\dd\mu(x)\dd\mu(y).
\end{align*}
Subtracting gives
\begin{align*}
&2\ip{d^k}{T_Uu}
-
2\int_\Omega d^ku^2\,\dd\mu
\\
&\qquad=
\int_{\Omega^2}
U(x,y)
\bigl(u(y)-u(x)\bigr)
\bigl(d(x)^k-d(y)^k\bigr)
\,\dd\mu(x)\dd\mu(y).
\end{align*}
Since
\[
d=1-u,
\]
the functions $u$ and $d^k=(1-u)^k$ are oppositely ordered.
Hence
\[
\bigl(u(y)-u(x)\bigr)
\bigl(d(x)^k-d(y)^k\bigr)\ge0.
\]
Thus \eqref{eq:correlation} follows.

Now
\[
u=1-d,
\]
so
\begin{align*}
\int_\Omega d^ku^2\,\dd\mu
&=
\int_\Omega d^k(1-d)^2\,\dd\mu
\\
&=
M_k-2M_{k+1}+M_{k+2}.
\end{align*}
Substituting into \eqref{eq:Jk-decomp} gives
\begin{align*}
J_k
&\ge
M_{k+1}-qM_k
+
M_k-2M_{k+1}+M_{k+2}
\\
&=
M_{k+2}-M_{k+1}+pM_k.
\end{align*}
Therefore \eqref{eq:Ik-Jk} yields
\[
I_k
\ge
2M_{k+2}-2M_{k+1}+pM_k.
\]
Finally,
\begin{align*}
&\left(
2M_{k+2}-2M_{k+1}+pM_k
\right)
-
\frac{2p-1}{p}M_{k+2}
\\
&\qquad=
\frac1p
\left(
M_{k+2}-2pM_{k+1}+p^2M_k
\right)
\\
&\qquad=
\frac1p
\int_\Omega
d(x)^k(d(x)-p)^2\,\dd\mu(x)
\\
&\qquad\ge0.
\end{align*}
This proves the asserted inequality.
\end{proof}

The previous lemma has the following useful interpretation.

\begin{corollary}[Average edge density of the links]
\label{cor:link-average}
Let $n\ge2$, and define
\[
M_n:=\int_\Omega d(x)^n\,\dd\mu(x).
\]
If $p>0$, then
\[
\frac{1}{M_n}
\int_\Omega d(x)^n z_x\,\dd\mu(x)
\ge
2-\frac1p.
\]
\end{corollary}

\begin{proof}
Since
\[
d(x)^nz_x=d(x)^{n-2}\tau(x),
\]
Lemma~\ref{lem:weighted-local-goodman}, applied with $k=n-2$,
gives
\[
\int d^nz_x
\ge
\frac{2p-1}{p}M_n.
\]
Since
\[
\frac{2p-1}{p}
=
2-\frac1p,
\]
the result follows.
\end{proof}


%%%%%%%%%%%%%%%%%%%%%%%%%%%%%%%%%%%%%%%%%%%%%%%%%%%%%%%%%%%%%%%%%%%%%%%%%%%%%%%
\section{A tangent-line extension lemma}
%%%%%%%%%%%%%%%%%%%%%%%%%%%%%%%%%%%%%%%%%%%%%%%%%%%%%%%%%%%%%%%%%%%%%%%%%%%%%%%

The next elementary observation allows us to use a lower bound that is
known only above a threshold $a$ even when some individual links have
edge density below $a$.

\begin{lemma}[Global tangent lower bound]
\label{lem:tangent}
Let $F$ be a finite graph and let
\[
\phi:[a,1]\to\R
\]
be differentiable and convex, where $0\le a<1$.  Assume that
\[
\phi(a)=0,
\qquad
\phi'(z)\ge0
\quad\text{for }z\in[a,1],
\]
and suppose that every graphon $V$ of edge density $z\ge a$ satisfies
\[
t(F,V)\ge\phi(z).
\]

Fix $c\in[a,1]$.  Then every graphon $V$, with arbitrary edge density
$z\in[0,1]$, satisfies
\[
\boxed{
t(F,V)
\ge
\phi(c)+\phi'(c)(z-c).
}
\]
\end{lemma}

\begin{proof}
Define the tangent line at $c$ by
\[
L_c(z)
=
\phi(c)+\phi'(c)(z-c).
\]

First suppose $z\ge a$.  Convexity of $\phi$ gives
\[
\phi(z)\ge L_c(z).
\]
Hence
\[
t(F,V)\ge\phi(z)\ge L_c(z).
\]

Now suppose $z<a$.  Again by convexity,
\[
L_c(a)\le\phi(a)=0.
\]
Since $\phi'(c)\ge0$ and $z<a$,
\[
L_c(z)\le L_c(a)\le0.
\]
On the other hand,
\[
t(F,V)\ge0.
\]
Thus
\[
t(F,V)\ge L_c(z)
\]
also in this case.
\end{proof}


%%%%%%%%%%%%%%%%%%%%%%%%%%%%%%%%%%%%%%%%%%%%%%%%%%%%%%%%%%%%%%%%%%%%%%%%%%%%%%%
\section{The universal-vertex lifting theorem}
%%%%%%%%%%%%%%%%%%%%%%%%%%%%%%%%%%%%%%%%%%%%%%%%%%%%%%%%%%%%%%%%%%%%%%%%%%%%%%%

We now prove the main abstract result.

\begin{theorem}[Universal-vertex lifting theorem]
\label{thm:universal-lifting}
Let $F$ be a graph on $n\ge2$ vertices, and write
\[
\phi(p):=\Phi_F(p).
\]
Suppose there exists $a\in[0,1)$ such that:

\begin{enumerate}[label=\textnormal{(\roman*)}]
\item
for every graphon $V$ with edge density $z\ge a$,
\[
t(F,V)\ge\phi(z);
\]

\item
for every $z\in[a,1]$,
\[
\phi(z)\ge0,
\qquad
\phi'(z)\ge0,
\qquad
\phi''(z)\ge0;
\]

\item
\[
\phi(a)=0.
\]
\end{enumerate}

Let
\[
F^+:=K_1\join F,
\qquad
b:=\frac1{2-a}.
\]
Then every graphon $W$ of edge density $p\ge b$ satisfies
\[
\boxed{
t(F^+,W)\ge\Phi_{F^+}(p).
}
\]
\end{theorem}

\begin{proof}
The case $p=1$ is immediate: then $W=1$ almost everywhere, hence
\[
t(F^+,W)=1=\Phi_{F^+}(1).
\]
Assume therefore
\[
b\le p<1.
\]

Define
\[
c:=2-\frac1p.
\]
Since
\[
b=\frac1{2-a},
\]
we have
\[
2-\frac1b=a.
\]
Moreover, the map
\[
p\longmapsto 2-\frac1p
\]
is increasing on $(0,1]$.  Hence
\[
p\ge b
\quad\Longrightarrow\quad
c\ge a.
\]
Clearly also $c<1$.

By Lemma~\ref{lem:link-identity},
\[
t(F^+,W)
=
\int_\Omega
d(x)^n t(F,W_x)\,\dd\mu(x).
\]

Apply Lemma~\ref{lem:tangent} to the graphon $W_x$, using the tangent
point $c$.  Its edge density is $z_x$, so
\[
t(F,W_x)
\ge
\phi(c)+\phi'(c)(z_x-c).
\]
Multiplying by $d(x)^n$ and integrating,
\begin{align*}
t(F^+,W)
&\ge
\phi(c)\int_\Omega d(x)^n\,\dd\mu(x)
\\
&\quad+
\phi'(c)
\left(
\int_\Omega d(x)^nz_x\,\dd\mu(x)
-
c\int_\Omega d(x)^n\,\dd\mu(x)
\right).
\end{align*}
Write
\[
M_n:=\int_\Omega d(x)^n\,\dd\mu(x).
\]
By Corollary~\ref{cor:link-average},
\[
\int_\Omega d(x)^nz_x\,\dd\mu(x)
\ge
\left(2-\frac1p\right)M_n
=
cM_n.
\]
Since $\phi'(c)\ge0$, the second term is nonnegative.  Therefore
\[
t(F^+,W)\ge M_n\phi(c).
\]

Since $x\mapsto x^n$ is convex on $[0,1]$, Jensen's inequality gives
\[
M_n
=
\int_\Omega d(x)^n\,\dd\mu(x)
\ge
\left(
\int_\Omega d(x)\,\dd\mu(x)
\right)^n
=
p^n.
\]
Because $\phi(c)\ge0$,
\[
t(F^+,W)\ge p^n\phi(c).
\]

It remains to identify $p^n\phi(c)$ with
$\Phi_{F^+}(p)$.

Put
\[
q:=1-p.
\]
From
\[
c=2-\frac1p
\]
we obtain
\[
1-c
=
\frac{q}{p},
\qquad
\frac1{1-c}
=
\frac{p}{q}.
\]
Hence
\[
\phi(c)
=
\left(\frac qp\right)^n
\chi_F\!\left(\frac pq\right),
\]
and therefore
\[
p^n\phi(c)
=
q^n\chi_F\!\left(\frac pq\right).
\]

On the other hand, if a universal vertex is added to $F$, then
\[
\chi_{F^+}(x)
=
x\,\chi_F(x-1).
\]
Indeed, the universal vertex may first be assigned any of the $x$
colours, and every vertex of $F$ must then be coloured using the
remaining $x-1$ colours.

Consequently,
\begin{align*}
\Phi_{F^+}(p)
&=
q^{n+1}
\chi_{F^+}\!\left(\frac1q\right)
\\
&=
q^{n+1}
\frac1q
\chi_F\!\left(\frac1q-1\right)
\\
&=
q^n
\chi_F\!\left(\frac pq\right)
\\
&=
p^n\phi(c).
\end{align*}
Thus
\[
t(F^+,W)\ge\Phi_{F^+}(p),
\]
as required.
\end{proof}


%%%%%%%%%%%%%%%%%%%%%%%%%%%%%%%%%%%%%%%%%%%%%%%%%%%%%%%%%%%%%%%%%%%%%%%%%%%%%%%
\section{Preservation of the convexity hypotheses}
%%%%%%%%%%%%%%%%%%%%%%%%%%%%%%%%%%%%%%%%%%%%%%%%%%%%%%%%%%%%%%%%%%%%%%%%%%%%%%%

The preceding theorem can be iterated because its analytic hypotheses
are themselves preserved.

\begin{proposition}[Convexity is preserved by adding a universal vertex]
\label{prop:preserve}
Assume the hypotheses of
Theorem~\ref{thm:universal-lifting}.
Let
\[
\psi(p):=\Phi_{F^+}(p),
\qquad
b=\frac1{2-a}.
\]
Then on $[b,1]$,
\[
\psi(p)\ge0,
\qquad
\psi'(p)\ge0,
\qquad
\psi''(p)\ge0,
\]
and
\[
\psi(b)=0.
\]
\end{proposition}

\begin{proof}
From the proof of
Theorem~\ref{thm:universal-lifting},
\[
\psi(p)
=
p^n
\phi\!\left(2-\frac1p\right).
\]
Set
\[
c(p):=2-\frac1p.
\]
For $p\in[b,1]$,
\[
c(p)\in[a,1].
\]
Thus $\phi(c(p))$, $\phi'(c(p))$, and $\phi''(c(p))$ are all
nonnegative.

Since
\[
c'(p)=\frac1{p^2},
\]
differentiating gives
\begin{align*}
\psi'(p)
&=
np^{n-1}\phi(c)
+
p^n\phi'(c)\frac1{p^2}
\\
&=
p^{n-2}
\left(
np\phi(c)+\phi'(c)
\right)
\ge0.
\end{align*}

Differentiating once more,
\begin{align*}
\psi''(p)
=
p^{n-4}
\bigl(
n(n-1)p^2\phi(c)
+
2(n-1)p\phi'(c)
+
\phi''(c)
\bigr).
\end{align*}
Every term in the bracket is nonnegative, so
\[
\psi''(p)\ge0.
\]
The nonnegativity of $\psi$ is immediate from
\[
\psi(p)=p^n\phi(c(p)).
\]

Finally,
\[
c(b)
=
2-\frac1b
=
a,
\]
so
\[
\psi(b)
=
b^n\phi(a)
=
0.
\]
\end{proof}


%%%%%%%%%%%%%%%%%%%%%%%%%%%%%%%%%%%%%%%%%%%%%%%%%%%%%%%%%%%%%%%%%%%%%%%%%%%%%%%
\section{Iteration}
%%%%%%%%%%%%%%%%%%%%%%%%%%%%%%%%%%%%%%%%%%%%%%%%%%%%%%%%%%%%%%%%%%%%%%%%%%%%%%%

The threshold transformation in the lifting theorem has a particularly
simple form.

\begin{lemma}[Threshold iteration]
\label{lem:threshold}
Let
\[
a_0=a,
\qquad
a_{j+1}:=\frac1{2-a_j}.
\]
Then
\[
\frac1{1-a_{j+1}}
=
\frac1{1-a_j}+1.
\]
Consequently,
\[
a_j
=
1-
\frac1{\frac1{1-a}+j}.
\]
In particular, if
\[
a=1-\frac1r,
\]
then
\[
a_j=1-\frac1{r+j}.
\]
\end{lemma}

\begin{proof}
Directly,
\begin{align*}
1-a_{j+1}
&=
1-\frac1{2-a_j}
\\
&=
\frac{1-a_j}{2-a_j}.
\end{align*}
Therefore
\[
\frac1{1-a_{j+1}}
=
\frac{2-a_j}{1-a_j}
=
1+\frac1{1-a_j}.
\]
Iteration gives the formula.
\end{proof}

Combining
Theorem~\ref{thm:universal-lifting},
Proposition~\ref{prop:preserve},
and Lemma~\ref{lem:threshold}
gives the following general statement.

\begin{corollary}[Iterated clique-join lifting]
\label{cor:iterated-lift}
Let $F$ satisfy the hypotheses of
Theorem~\ref{thm:universal-lifting}
on $[a,1]$.
For $s\ge0$, define
\[
F_s:=K_s\join F.
\]
Then
\[
t(F_s,W)\ge\Phi_{F_s}(p)
\]
for every graphon $W$ whose edge density satisfies
\[
p\ge
1-
\frac1{\frac1{1-a}+s}.
\]
\end{corollary}

\begin{proof}
Set $F_0=F$.  Since
\[
F_{j+1}
=
K_1\join F_j,
\]
Theorem~\ref{thm:universal-lifting} gives the desired inequality at
the next threshold, while
Proposition~\ref{prop:preserve} shows that all hypotheses required
for the next application remain valid.
Lemma~\ref{lem:threshold} gives the stated closed form for the
threshold.
\end{proof}


%%%%%%%%%%%%%%%%%%%%%%%%%%%%%%%%%%%%%%%%%%%%%%%%%%%%%%%%%%%%%%%%%%%%%%%%%%%%%%%
\section{Convexity of the odd-cycle target}
%%%%%%%%%%%%%%%%%%%%%%%%%%%%%%%%%%%%%%%%%%%%%%%%%%%%%%%%%%%%%%%%%%%%%%%%%%%%%%%

We now verify that odd cycles satisfy the analytic assumptions needed
for the lifting theorem.

\begin{lemma}[Convexity of the odd-cycle Goodman polynomial]
\label{lem:odd-convexity}
Let $m\ge3$ be odd and define
\[
g_m(p)
=
p^m-p(1-p)^{m-1}.
\]
Then on $[1/2,1]$,
\[
g_m(p)\ge0,
\qquad
g_m'(p)>0,
\qquad
g_m''(p)>0,
\]
and
\[
g_m(1/2)=0.
\]
\end{lemma}

\begin{proof}
Write
\[
q:=1-p.
\]
If $p\ge1/2$, then $p\ge q$, and therefore
\[
g_m(p)
=
p\left(p^{m-1}-q^{m-1}\right)
\ge0.
\]
At $p=1/2$,
\[
g_m(1/2)=0.
\]

Differentiate:
\begin{align*}
g_m'(p)
&=
mp^{m-1}
-
q^{m-1}
+
(m-1)pq^{m-2},
\end{align*}
and
\begin{align}
g_m''(p)
&=
(m-1)
\left(
mp^{m-2}
+
2q^{m-2}
-
(m-2)pq^{m-3}
\right).
\tag{7.1}
\label{eq:g-second}
\end{align}

Suppose first that
\[
\frac12\le p<1.
\]
Set
\[
R:=\frac pq\ge1.
\]
Factoring $q^{m-2}$ from
\eqref{eq:g-second} gives
\[
g_m''(p)
=
(m-1)q^{m-2}
\left(
mR^{m-2}
+
2
-
(m-2)R
\right).
\]
Since $m\ge3$ and $R\ge1$,
\[
R^{m-2}\ge R.
\]
Hence
\begin{align*}
mR^{m-2}+2-(m-2)R
&\ge
mR+2-(m-2)R
\\
&=
2R+2
\\
&>0.
\end{align*}
Thus
\[
g_m''(p)>0
\]
for $1/2\le p<1$.  The endpoint $p=1$ follows by continuity, or
directly from \eqref{eq:g-second}.

Therefore $g_m'$ is increasing on $[1/2,1]$.  At the left endpoint,
\begin{align*}
g_m'\!\left(\frac12\right)
&=
\frac{m}{2^{m-1}}
-
\frac1{2^{m-1}}
+
\frac{m-1}{2^{m-1}}
\\
&=
\frac{m-1}{2^{m-2}}
\\
&>0.
\end{align*}
Hence
\[
g_m'(p)>0
\]
throughout $[1/2,1]$.
\end{proof}


%%%%%%%%%%%%%%%%%%%%%%%%%%%%%%%%%%%%%%%%%%%%%%%%%%%%%%%%%%%%%%%%%%%%%%%%%%%%%%%
\section{Clique joins of odd cycles}
%%%%%%%%%%%%%%%%%%%%%%%%%%%%%%%%%%%%%%%%%%%%%%%%%%%%%%%%%%%%%%%%%%%%%%%%%%%%%%%

We can now state and prove the main concrete theorem.

\begin{theorem}[Clique joins of odd cycles]
\label{thm:clique-join-cycle}
Let $m\ge3$ be odd and let $s\ge0$ be an integer.  Define
\[
H_{s,m}:=K_s\join C_m.
\]
Then every graphon $W$ with
\[
p=t(K_2,W)
\ge
1-\frac1{s+2}
\]
satisfies
\[
\boxed{
t(H_{s,m},W)
\ge
\Phi_{H_{s,m}}(p)
=
(1-p)^{m+s}
\chi_{H_{s,m}}
\!\left(\frac1{1-p}\right).
}
\]

Moreover,
\[
\chi(H_{s,m})=s+3.
\]
Consequently the same inequality certainly holds throughout the
relaxed range
\[
p\ge
1-\frac1{\chi(H_{s,m})}
=
1-\frac1{s+3}.
\]
\end{theorem}

\begin{proof}
Take
\[
F=C_m.
\]
By Theorem~\ref{thm:odd-cycle-input},
\[
t(C_m,W)\ge g_m(p)=\Phi_{C_m}(p)
\]
for every graphon $W$.

By Lemma~\ref{lem:odd-convexity}, on
\[
\left[\frac12,1\right]
\]
the function
\[
\phi(p):=\Phi_{C_m}(p)=g_m(p)
\]
satisfies
\[
\phi\ge0,
\qquad
\phi'\ge0,
\qquad
\phi''\ge0,
\qquad
\phi(1/2)=0.
\]
Thus all hypotheses of
Corollary~\ref{cor:iterated-lift}
hold with
\[
a=\frac12.
\]

Since
\[
\frac1{1-a}=2,
\]
after adjoining $s$ mutually adjacent universal vertices the threshold
becomes
\[
1-\frac1{2+s}.
\]
The resulting graph is exactly
\[
K_s\join C_m.
\]
Therefore
\[
t(K_s\join C_m,W)
\ge
\Phi_{K_s\join C_m}(p)
\]
whenever
\[
p\ge1-\frac1{s+2}.
\]

It remains only to determine the chromatic number.
Every vertex of $K_s$ is adjacent to every vertex of $C_m$.
Thus the colours used on $K_s$ and on $C_m$ must be disjoint.
Since
\[
\chi(K_s)=s
\]
and, for odd $m$,
\[
\chi(C_m)=3,
\]
we obtain
\[
\chi(K_s\join C_m)=s+3.
\]
Since
\[
1-\frac1{s+3}
>
1-\frac1{s+2},
\]
the requested relaxed density regime is contained in the range already
proved.
\end{proof}


%%%%%%%%%%%%%%%%%%%%%%%%%%%%%%%%%%%%%%%%%%%%%%%%%%%%%%%%%%%%%%%%%%%%%%%%%%%%%%%
\section{Explicit form of the lower bound}
%%%%%%%%%%%%%%%%%%%%%%%%%%%%%%%%%%%%%%%%%%%%%%%%%%%%%%%%%%%%%%%%%%%%%%%%%%%%%%%

The chromatic polynomial of the graph in
Theorem~\ref{thm:clique-join-cycle}
can be written explicitly.

For $s\ge0$, define
\[
A_s(p)
:=
\prod_{j=0}^{s-1}
\bigl(1-j(1-p)\bigr),
\]
where the empty product $A_0(p)$ is $1$.

Also define
\[
q:=1-p
\]
and
\[
\rho_s(p)
:=
1-(s+1)q
=
(s+1)p-s.
\]

\begin{proposition}[Explicit target polynomial]
\label{prop:explicit}
For odd $m\ge3$,
\[
\boxed{
\Phi_{K_s\join C_m}(p)
=
A_s(p)
\left(
\rho_s(p)^m
-
\rho_s(p)q^{m-1}
\right).
}
\]
Consequently,
\[
\boxed{
t(K_s\join C_m,W)
\ge
A_s(p)
\left(
\rho_s(p)^m
-
\rho_s(p)(1-p)^{m-1}
\right)
}
\]
whenever
\[
p\ge1-\frac1{s+2}.
\]
\end{proposition}

\begin{proof}
For arbitrary graphs $G,H$,
\[
\chi_{G\join H}(x)
\]
is obtained by choosing disjoint palettes for $G$ and $H$.
In the special case $G=K_s$,
\[
\chi_{K_s\join H}(x)
=
(x)_s\chi_H(x-s),
\]
where
\[
(x)_s=x(x-1)\cdots(x-s+1).
\]
Hence
\[
\chi_{K_s\join C_m}(x)
=
(x)_s\chi_{C_m}(x-s).
\]

The chromatic polynomial of a cycle is
\[
\chi_{C_m}(y)
=
(y-1)^m+(-1)^m(y-1).
\]
Since $m$ is odd,
\[
\chi_{C_m}(y)
=
(y-1)^m-(y-1).
\]
Therefore
\[
\chi_{K_s\join C_m}(x)
=
(x)_s
\left(
(x-s-1)^m-(x-s-1)
\right).
\]

Now put
\[
x=\frac1q,
\qquad
q=1-p.
\]
Then
\[
q^s\left(\frac1q\right)_s
=
\prod_{j=0}^{s-1}(1-jq)
=
A_s(p).
\]
Furthermore,
\[
q\left(\frac1q-s-1\right)
=
1-(s+1)q
=
\rho_s(p).
\]
Thus
\begin{align*}
\Phi_{K_s\join C_m}(p)
&=
q^{m+s}
\chi_{K_s\join C_m}\!\left(\frac1q\right)
\\
&=
A_s(p)
q^m
\left[
\left(\frac{\rho_s(p)}q\right)^m
-
\frac{\rho_s(p)}q
\right]
\\
&=
A_s(p)
\left(
\rho_s(p)^m
-
\rho_s(p)q^{m-1}
\right).
\end{align*}
This is the claimed formula.
\end{proof}


%%%%%%%%%%%%%%%%%%%%%%%%%%%%%%%%%%%%%%%%%%%%%%%%%%%%%%%%%%%%%%%%%%%%%%%%%%%%%%%
\section{Some special cases}
%%%%%%%%%%%%%%%%%%%%%%%%%%%%%%%%%%%%%%%%%%%%%%%%%%%%%%%%%%%%%%%%%%%%%%%%%%%%%%%

\subsection{Odd wheels}

Taking $s=1$, we obtain the join of a universal vertex with an odd
cycle.  Since
\[
A_1(p)=1,
\qquad
\rho_1(p)=2p-1,
\]
Theorem~\ref{thm:clique-join-cycle} gives
\[
t(K_1\join C_m,W)
\ge
(2p-1)^m
-
(2p-1)(1-p)^{m-1}
\]
for every odd $m\ge3$ and every
\[
p\ge\frac23.
\]
For $m\ge5$ these are nonchordal $4$-chromatic graphs.

\subsection{Two universal vertices}

Taking $s=2$ gives
\[
A_2(p)=p,
\qquad
\rho_2(p)=3p-2.
\]
Therefore
\[
t(K_2\join C_m,W)
\ge
p\left(
(3p-2)^m
-
(3p-2)(1-p)^{m-1}
\right)
\]
for
\[
p\ge\frac34.
\]
These graphs have chromatic number $5$.

\subsection{Arbitrary chromatic number}

For every $r\ge3$, take
\[
s=r-3.
\]
Then
\[
\chi(K_{r-3}\join C_m)=r,
\]
and
\[
t(K_{r-3}\join C_m,W)
\ge
\Phi_{K_{r-3}\join C_m}(p)
\]
for every
\[
p\ge
1-\frac1{r-1}.
\]
In particular, the conclusion holds in the weaker requested regime
\[
p\ge1-\frac1r.
\]

For $m\ge5$, the induced subgraph on the cycle vertices is $C_m$.
Hence $K_{r-3}\join C_m$ is not chordal.
Thus these examples are not consequences of the pure-chordal theorem.


%%%%%%%%%%%%%%%%%%%%%%%%%%%%%%%%%%%%%%%%%%%%%%%%%%%%%%%%%%%%%%%%%%%%%%%%%%%%%%%
\section{Sharpness at the chromatic Tur\'an point}
%%%%%%%%%%%%%%%%%%%%%%%%%%%%%%%%%%%%%%%%%%%%%%%%%%%%%%%%%%%%%%%%%%%%%%%%%%%%%%%

Let
\[
r=s+3=\chi(K_s\join C_m).
\]
Consider the balanced complete $r$-partite graphon $T_r$.
Its edge density is
\[
p_r=1-\frac1r.
\]
For every graph $H$,
\[
t(H,T_r)
=
\frac{\chi_H(r)}{r^{v(H)}}.
\]
Since
\[
1-p_r=\frac1r,
\]
we have
\[
t(H,T_r)
=
(1-p_r)^{v(H)}
\chi_H\!\left(\frac1{1-p_r}\right)
=
\Phi_H(p_r).
\]
Applying this to
\[
H=K_s\join C_m
\]
shows that the bound of
Theorem~\ref{thm:clique-join-cycle}
is attained at
\[
p=1-\frac1{s+3}.
\]
Thus the theorem contains, in particular, the natural equality point
corresponding to the relaxed threshold in the motivating question.


%%%%%%%%%%%%%%%%%%%%%%%%%%%%%%%%%%%%%%%%%%%%%%%%%%%%%%%%%%%%%%%%%%%%%%%%%%%%%%%
\section{Summary of the mechanism}
%%%%%%%%%%%%%%%%%%%%%%%%%%%%%%%%%%%%%%%%%%%%%%%%%%%%%%%%%%%%%%%%%%%%%%%%%%%%%%%

The proof consists of four exact ingredients.

\begin{enumerate}
\item
For $F^+=K_1\join F$, conditioning on the image $x$ of the new universal
vertex yields
\[
t(F^+,W)
=
\int d(x)^{v(F)}t(F,W_x)\,\dd\mu(x).
\]

\item
The weighted rooted-triangle inequality gives
\[
\frac{
\int d(x)^{v(F)}t(K_2,W_x)\,\dd\mu(x)
}{
\int d(x)^{v(F)}\,\dd\mu(x)
}
\ge
2-\frac1p.
\]

\item
If $\Phi_F$ is increasing and convex above a zero $a$, then the tangent
line at any $c\ge a$ lies below $t(F,V)$ for \emph{all} graphons $V$,
including those whose edge density is below $a$.

\item
The transformation
\[
c=2-\frac1p
\]
satisfies
\[
p^{v(F)}\Phi_F(c)
=
\Phi_{K_1\join F}(p).
\]
The corresponding threshold transforms as
\[
a\longmapsto \frac1{2-a}.
\]
If
\[
a=1-\frac1r,
\]
this becomes
\[
1-\frac1r
\longmapsto
1-\frac1{r+1}.
\]
\end{enumerate}

For odd cycles the initial zero is $a=1/2$, and the odd-cycle target
is increasing and convex on $[1/2,1]$.  Iterating the construction
therefore yields the entire family
\[
K_s\join C_{2k+1}.
\]

%%%%%%%%%%%%%%%%%%%%%%%%%%%%%%%%%%%%%%%%%%%%%%%%%%%%%%%%%%%%%%%%%%%%%%%%%%%%%%%
\appendix
\section{What the Lean formalization actually proves}
%%%%%%%%%%%%%%%%%%%%%%%%%%%%%%%%%%%%%%%%%%%%%%%%%%%%%%%%%%%%%%%%%%%%%%%%%%%%%%%

What the six-vertex catalogue takes from this note is the single case
$(s,m)=(1,5)$, that is $K_1\join C_5$, which is Atlas $187$ and is
\texttt{verified}.  The development is
\texttt{lean/Taeyoung/Methods/OddCycleCone.lean} over
\texttt{Methods/OddCycleC5/}.  Two differences of substance.

\subsection*{Theorem~\ref{thm:odd-cycle-input} is not assumed, and not proved
either}

The standing input is used at $m=5$ only, and there it is not taken as an
input: $t(C_5,W)\geq p^5-p(1-p)^4$ is proved outright in
\texttt{Methods/OddCycleC5/}, from integral kernel definitions, through a
necklace identity, a sum-of-squares certificate for the length-five moments,
and edge deletion.  That development is deliberately restricted to length five
and invokes no theorem about all odd cycles, so $g_m$ for $m\geq7$ is neither
assumed nor available.

\subsection*{Theorem~\ref{thm:universal-lifting} is not formalized in general}

Neither the universal lifting theorem, nor its iteration
(Corollary~\ref{cor:iterated-lift}), nor Theorem~\ref{thm:clique-join-cycle}
for general $(s,m)$, is stated in Lean.  Its analytic content --- the weighted
local Goodman inequality of Lemma~\ref{lem:weighted-local-goodman}, the link
identity, and the link-average corollary --- is present, and is shared with
three other cone families through the abstract chain of
\texttt{Methods/ConeBound.lean}; what is instantiated from it here is one row.

That instantiation also avoids Lemma~\ref{lem:odd-convexity}.  No convexity of
$\Phi_{C_m}$ is verified.  Instead, at $m=5$,
\[
 \Phi_{C_5}(z)=z^5-z(1-z)^4=z(2z-1)(2z^2-2z+1)
\]
is a quartic, and the remainder of its tangent at $c$ factors as
\[
 \Phi_{C_5}(w)-\Phi_{C_5}(c)-\Phi_{C_5}'(c)(w-c)
 =(w-c)^2\bigl(4w^2+8wc+12c^2-6w-12c+4\bigr),
\]
whose second factor, rewritten in $w-\tfrac12$ and $c-\tfrac12$, is
\[
 1+2\bigl(w-\tfrac12\bigr)+4\bigl(c-\tfrac12\bigr)
 +4\bigl(w-\tfrac12\bigr)^2+8\bigl(w-\tfrac12\bigr)\bigl(c-\tfrac12\bigr)
 +12\bigl(c-\tfrac12\bigr)^2
\]
with every coefficient nonnegative.  The affine minorant on $[1/2,1]$ follows
by inspection.  Note also that $2z^2-2z+1=z^2+(1-z)^2>0$, so $\Phi_{C_5}\geq0$
exactly on $[1/2,1]$.

\end{document}